\documentclass[11pt,a4paper]{letter}
\usepackage[utf8]{inputenc}
\usepackage[T1]{fontenc}
\usepackage{mlmodern}
\usepackage{amsmath}
\usepackage{amsfonts}
\usepackage{amssymb}
\usepackage{graphicx}
\usepackage{geometry}
\usepackage{csquotes}
\usepackage[final]{microtype}
\geometry{margin=1in}

\address{Aniara Lyn Lauron \\ Gåseholmvej 61 \\ Herlev, Denmark \\ Email: aniaralyn.lauron@gmail.com \\ Phone: +45 91 80 63 94}
\date{February 22, 2026}

\begin{document}

\begin{letter}{Admissions Committee \\ University of Southern Denmark (SDU) \\ SDU Vejle Campus \\ Denmark}

\opening{Dear Admissions Committee,}

I am writing to express my strong interest in the Bachelor of Science in Artificial Intelligence programme at the University of Southern Denmark (SDU), starting in the fall of 2026. I hold a Bachelor of Science in Pharmacy from the University of Southern Philippines and am eager to transition into the field of AI, where I can combine my scientific foundation with data driven innovation and intelligent systems. The programme's description emphasises that programming experience is not required before starting, what matters most is an interest in learning and a strong aptitude for logic and mathematics. This matches my current mindset perfectly, as I am motivated to acquire coding skills and apply them to real problems.

My academic foundation in pharmacy has equipped me with a solid understanding of the scientific method, data analysis, and ethical responsibility. During my studies I conducted research on the anti-inflammatory effects of natural extracts, culminating in my thesis \enquote{The Effect of Hagonoy Leaves Extract on the Inflamed Paws of Male Albino Rabbits}. The project involved chemical extraction, animal experimentation, and statistical evaluation experiences that honed my analytical thinking and gave me an appreciation for modelling complex systems. Over time I became fascinated by the idea of teaching machines to learn from data, a concept that sits at the heart of artificial intelligence.

Beyond academics, I have gained practical experience in diverse roles. As a pharmacy technician at Optum Global Solutions Inc. and Collabera Technologies Private Limited I managed medication dispensing and collaborated across healthcare teams tasks that required precision, reliability and teamwork. My freelance work as an English teacher helped develop my communication skills, and my current role as a hotel assistant at The Ellen Group has enhanced my problem solving and customer service abilities. Each position taught me how to approach problems methodically and adapt to new tools, which I believe is essential for studying AI.

Being surrounded by engineers and technologists during my international experiences sparked my passion for technology and intelligent systems. Living in the Philippines, South Korea, the Netherlands, and Denmark has shaped my adaptability and global perspective. The Philippines fostered resilience and community values, Korea's dynamic tech culture inspired innovation, the Netherlands' focus on sustainability motivated my desire to apply computation to real problems, and Denmark's collaborative, pedagogy-driven approach aligns with how I learn best. I am driven by a genuine desire to help others and to continuously improve, qualities that will serve me well in the collaborative, project-based environment at SDU Vejle.

To bridge my pharmacy background with AI interests, I have pursued hands-on technical projects. In 2024 I designed and built an Arduino based automated plant irrigation system, programming firmware in C++ to automate watering schedules and integrating sensors and actuators. The project taught me the basics of input/output, control loops and embedded programming concepts that naturally extend to intelligent agents and robotics.

Additionally, my 2026 custom piano computer workstation project involves woodworking, design, and assembly, reflecting creativity and precision in crafting functional solutions. My 2023 silver ring making project demonstrated metalworking and engraving skills, highlighting my flexibility and capacity to learn new disciplines. These diverse hobbies reinforce my belief that intelligence, whether human or artificial, is built through experience and experimentation.

I believe that SDU Vejle's BSc programme in Artificial Intelligence is the ideal setting for my transition. The course emphasises theory, development, and application of AI methods topics such as optimisation, machine learning, logic, algorithms, and ethics resonate strongly with my interests. The international environment at the new IT campus, the close proximity to leading employers of IT specialists, and the accessible lecturers create a supportive atmosphere for motivated students. I am excited by the programme's focus on collaborative, project based work and the opportunity to work with companies on real world problems.

Looking ahead, upon completing the BSc in Artificial Intelligence, I plan to pursue a master's degree and eventually work at the intersection of AI and life sciences. My long term goal is to develop intelligent tools and diagnostic systems that leverage my pharmacy background to improve healthcare and quality of life. I envision contributing to Denmark's thriving technology sector, perhaps at companies like Novo Nordisk or Coloplast, by building solutions that combine data science, machine learning, and domain expertise. Graduates of the programme enjoy strong job security and career opportunities both in Denmark and abroad, and I am eager to join that community.

I am excited about the opportunity to join SDU and contribute my unique perspective to the Artificial Intelligence programme. Thank you for considering my application. I look forward to the possibility of discussing how my background and aspirations align with SDU's mission.

\closing{Sincerely, \\
Aniara Lyn Lauron}

\end{letter}

\end{document}